\section{The Tower}

This appendix records the doubling construction used throughout, the single move that builds the whole algebra tower and that the pinch-point argument rests on.

\subsection*{The doubling}

Given an algebra $A$ with a conjugation $x \mapsto \bar{x}$, the Cayley-Dickson double is the set of pairs $(a, b)$ with $a, b$ in $A$, with multiplication

\[ (a, b)(c, d) = (ac - \bar{d}\,b,\ da + b\bar{c}) \]

and conjugation $\overline{(a, b)} = (\bar{a}, -b)$. The double has twice the dimension of $A$. Applied repeatedly from the one-dimensional start it gives the tower

\[ \mathbb{Z} \to \mathbb{C} \to \mathbb{H} \to \mathbb{O} \to \mathbb{S}, \qquad 1, 2, 4, 8, 16. \]

the integers, complexes, quaternions, octonions, and sedenions. At the base the reals are read discretely as the integers, since the construction is discrete throughout and the continuum is emergent.

\subsection*{What each step loses}

Each doubling trades away one property, in a fixed order.

\begin{center}
\begin{tabular}{lll}
Step & Dimension & Property lost \\
\hline
Z to C & 1 to 2 & order (gains a square root of minus one) \\
C to H & 2 to 4 & commutativity \\
H to O & 4 to 8 & associativity \\
O to S & 8 to 16 & the norm (composition) \\
\end{tabular}
\end{center}

The order is forced by one fact. The double of $A$ keeps the norm-multiplicative property if and only if $A$ is associative, the Hurwitz condition. The integers, complexes, and quaternions are associative, so the first three doublings keep the norm, producing the complexes, quaternions, and octonions. The octonions are not associative, so the fourth doubling loses the norm, and the sedenions have zero divisors, with the explicit pair

\[ (e_1 + e_{10})(e_5 + e_{14}) = 0. \]

a product of two nonzero sedenions that vanishes. The composition algebras are therefore exactly the integers, complexes, quaternions, and octonions, of dimensions one, two, four, eight. This is the ceiling at eight.

\subsection*{The floor and the pinch}

Reversibility, the demand that distinctions fold without loss, is exactly the norm-multiplicative condition, so it caps the dimension at eight. Carrying three mirrored families of spin needs the threefold symmetry called triality, which lives only in the eight-dimensional rotation structure, so three-family matter floors the dimension at eight. The ceiling and the floor meet at the octonions, the last reversible rung before zero divisors appear and reversibility shatters. The octonions are therefore not a choice but a pinch-point.

\subsection*{The octonion multiplication table}

For reference, the multiplication of the seven imaginary units used throughout the fermion construction. Each unit squares to minus one, and the off-diagonal products follow the Fano-plane orientation. The entry in row $e_i$, column $e_j$ is the product $e_i e_j$.

\begin{center}
\begin{tabular}{c|ccccccc}
 & $e_1$ & $e_2$ & $e_3$ & $e_4$ & $e_5$ & $e_6$ & $e_7$ \\
\hline
$e_1$ & $-1$ & $e_3$ & $-e_2$ & $e_5$ & $-e_4$ & $-e_7$ & $e_6$ \\
$e_2$ & $-e_3$ & $-1$ & $e_1$ & $e_6$ & $e_7$ & $-e_4$ & $-e_5$ \\
$e_3$ & $e_2$ & $-e_1$ & $-1$ & $e_7$ & $-e_6$ & $e_5$ & $-e_4$ \\
$e_4$ & $-e_5$ & $-e_6$ & $-e_7$ & $-1$ & $e_1$ & $e_2$ & $e_3$ \\
$e_5$ & $e_4$ & $-e_7$ & $e_6$ & $-e_1$ & $-1$ & $-e_3$ & $e_2$ \\
$e_6$ & $e_7$ & $e_4$ & $-e_5$ & $-e_2$ & $e_3$ & $-1$ & $-e_1$ \\
$e_7$ & $-e_6$ & $e_5$ & $e_4$ & $-e_3$ & $-e_2$ & $e_1$ & $-1$ \\
\end{tabular}
\end{center}

The table is antisymmetric off the diagonal, recording non-commutativity, and a direct check of any associator, such as $(e_1 e_2) e_4$ against $e_1 (e_2 e_4)$, shows non-associativity. These two facts, the anticommutation and the non-associativity, are the inputs to the Clifford and triality arguments. The seven left-multiplication matrices that generate the Clifford algebra of the fermion section are read column by column from this table.
